\documentclass[11pt]{article}

\usepackage[margin=0.6in]{geometry}
\usepackage[T1]{fontenc}
\usepackage[utf8]{inputenc}
\usepackage{lmodern}
\usepackage{array}
\usepackage{hyperref}
\usepackage{parskip}
\usepackage{ragged2e}
\usepackage{setspace}
\usepackage[table]{xcolor}
\usepackage[natbibapa]{apacite}

\definecolor{PrimaryRed}{HTML}{8B1E1E}
\definecolor{SoftBlue}{HTML}{DDEEFF}

\hypersetup{colorlinks=true,linkcolor=black,urlcolor=black,citecolor=black}
\setlength{\parindent}{0pt}
\setstretch{1.08}
\bibliographystyle{apacite}
\newcommand{\reviewsection}[1]{\vspace{0.07in}\noindent{\normalsize\bfseries\color{PrimaryRed}#1}\par\vspace{0.02in}}

\begin{document}

\begin{titlepage}
\newgeometry{margin=0.7in}
\thispagestyle{empty}

\begin{center}
{\large University of Rochester\\}
{\normalsize Warner School of Education\\[0.15in]}

{\color{PrimaryRed}\rule{\textwidth}{1.4pt}}\\[0.12in]
{\huge\bfseries Final Course Reflection\\[0.08in]
\Huge Inclusive Instructional Design,\\[0.08in]
Causal Agency, and Evidence That Sees Students}\\[0.18in]
{\Large A Personal and Professional Reflection Jornal on ED452C}\\[0.12in]

\vfill

{\color{PrimaryRed}\rule{0.78\textwidth}{0.7pt}}\\[0.22in]

\colorbox{SoftBlue}{%
\parbox{0.88\textwidth}{%
\centering\small\vspace{0.02in}
\textbf{From Individual Accommodation to Environment-Level Design}\\[0.05in]
Reflecting on how ED452C reshaped my understanding of inclusion through Universal Design for Learning, strengths-based practice, self-determination, social-emotional learning, inclusive assessment, transition planning, and my own Puzzle Plan and EQUITAS lenses for neurodivergent, culturally and linguistically diverse, disabled, chronically ill, multilingual, and multiply identified students.
}%
}
\end{center}

\vfill

\begin{center}
\renewcommand{\arraystretch}{1.18}
\setlength{\tabcolsep}{4pt}
\begin{tabular}{>{\bfseries\RaggedLeft}p{1.65in} p{4.9in}}
Student & Piter Zacari Garcia Bautista \\
Assignment & Final Course Reflection \\
Course & ED452C --- Instructional Strategies for Inclusive Classrooms C \\
Focus & Inclusive instructional design as access, participation, growth, belonging, and agency \\
Lens & Puzzle Plan and EQUITAS; equity-centered, neurodiversity-affirming, trauma-responsive practice \\
Submission Date & June 2026 \\
\end{tabular}
\end{center}

\vfill

\begin{center}
\begin{minipage}{0.92\textwidth}
\centering\footnotesize
\itshape
Inclusion is not proven by where a student sits. It is proven by what the environment allows the student to access, decide, express, revise, and become.\\
{\color{PrimaryRed}\rule{4in}{0.8pt}}
\end{minipage}
\end{center}

\vfill

\begin{center}
\begin{minipage}{1.0\textwidth}
\small
This final reflection answers the three required ED452C questions: what course topics I now understand well, what topics I want to learn more about, and how I will pursue further learning. I connect those questions to the course readings, my Pine Brook IGNITE teaching-placement patterns, my Learning Assistant lesson on evidence in support of inclusive education, and the journal-review theme of classroom climate as an environment-level intervention.
\end{minipage}
\end{center}

\restoregeometry
\end{titlepage}
\clearpage

\reviewsection{Topics I Now Understand Well}

The topic I understand most strongly after ED452C is that inclusive education is not a place where students are sent; it is a design responsibility. Wehmeyer and Kurth's framing of special education helped me separate inclusion from simple physical placement. A student can sit in a general education classroom and still be excluded by how directions are delivered, how quickly students are expected to respond, how evidence is collected, or which forms of participation are treated as legitimate. That shift changed my first professional question. Instead of asking, ``What is wrong with this student?'' I now ask, ``What is this environment demanding, and whose thinking is being hidden by that demand?'' I also developed a stronger understanding of disability through a person-environment fit lens. The social-ecological view makes it harder to blame the learner when the real barrier may be speed, language load, sensory demand, executive-function demand, unfamiliar technology, or narrow definitions of competence. In my Pine Brook IGNITE placement, I saw this repeatedly. During hands-on STEM lessons, students who hesitated were not always confused about the content. Some were decoding multi-step verbal directions, negotiating group roles, handling unfamiliar materials, or waiting for a visual cue that made the task make sense. In the AI project brainstorming lesson, ``stuck'' looked similar on the surface even when the barrier was different: a login issue, too many AI-generated choices, uncertainty about connecting a personal interest to computer science, or simple overwhelm. ED452C gave me language for what I was noticing: behavior alone does not tell the whole story.

A second topic I now understand well is self-determination as causal agency. Wehmeyer and Shogren helped me see self-determination not as a personality trait or a transition-planning milestone, but as something built through repeated experiences of choosing, acting, evaluating, and adjusting. In the AI brainstorming lesson, one student who had been quiet during verbal setup became more engaged once the task began with a personal-interest choice. He chose Minecraft, read the AI output, judged that the first ideas were too complex, and asked the AI to simplify them without being told. The key design move was not the technology itself. It was that the student's preference was the entry point. He did not have to earn the right to decide. This matters especially for neurodivergent learners, culturally and linguistically diverse students, multilingual students, disabled students, and students navigating multiple identities. These students are often over-scaffolded or over-directed in the name of support. ED452C helped me name the difference between scaffolding that opens participation and scaffolding that quietly takes agency away. Support should reduce unnecessary barriers, but it should not remove the student's opportunity to practice judgment, voice, and authorship.

The third topic I understand more deeply is evidence in inclusive education. My Learning Assistant lesson helped me synthesize this part of the course through the frame of access, participation, growth, and belonging. I no longer think task completion alone is enough. In the Lego sail car lesson at Pine Brook, the build phase took the full period, and the redesign phase---the part most connected to student decision-making, testing, revision, and agency---was the first thing lost to time. That helped me understand that evidence is also a design choice. If a lesson only measures who finished fastest, it may miss who revised a strategy, explained an idea to a peer, asked a better question, or needed one different access point to show what they knew.

My Journal Article Review strengthened this same understanding through the RULER study. Rivers and colleagues (2013) helped me see social-emotional learning as more than individual skill instruction. When implemented well, SEL can reshape classroom climate by increasing warmth, responsiveness, student autonomy, and belonging. That connects directly to my Puzzle Plan and EQUITAS lens: the environment is part of the intervention. For students with ADHD, chronic illness, trauma histories, anxiety, language differences, or multiple marginalized identities, classroom climate is not secondary to learning. It can determine whether learning feels possible.

\newpage
\reviewsection{Topics I Want to Learn More About}

The main topic I want to learn more about is how to audit lessons for agency, not only access. Before this course, I often thought first about whether all students could reach the task: Are the directions clear? Are the materials available? Is there a scaffold? Those questions still matter, but ED452C pushed me toward a harder question: Does every student get to be the one who decides? I want to become better at designing lessons where choice, revision, voice, and self-monitoring are not optional extensions that disappear when time runs short. They need to be built into the core structure of the lesson. I also want to learn more about the balance between scaffolding and over-scaffolding. This is especially important for students who are neurodivergent, multilingual, disabled, or labeled as struggling. Teachers often want to help, but help can become control if it prevents students from practicing judgment. A visual checklist, partner talk structure, or prompt frame can support access without taking over the student's thinking. But an adult who immediately tells the student what to do next may reduce frustration while also reducing agency. That tension is something I want to keep investigating. A third topic I want to keep studying is inclusive assessment, especially evidence plans that value process as well as product. I want to better understand how teachers can collect evidence across oral explanation, peer collaboration, revision, debugging talk, planning artifacts, questions, confidence, and belonging over time. This connects to my broader equity work because many students are misread when assessment privileges speed, written output, compliance, or standard academic English. I also want to continue learning how technology and AI can support inclusion without creating new inequities through login barriers, interface complexity, language demands, privacy concerns, or too many choices.

\reviewsection{How I Will Pursue Further Learning}

I will pursue further learning through practice, documentation, and research. First, I will keep using my teaching placements as a place to test inclusive design questions in real time. For each lesson I plan or support, I want to ask three audit questions before teaching: What are the likely access barriers? Where is the student's agency built into the task? What evidence will I collect beyond completion? These questions are simple enough to use consistently, but they connect directly to UDL, strengths-based design, self-determination, and inclusive assessment. Second, I will keep building portfolio-ready documentation from my teaching work. ED452C helped me see that every discussion post, lesson reflection, and placement note can become a small piece of a larger professional evidence system. My goal is to continue de-identifying classroom patterns and turning them into reusable teaching evidence. That means recording not just what happened, but what the moment taught me about access, agency, participation, and belonging. This will help me build a stronger certification portfolio while also helping me grow as a reflective teacher. Third, I will connect ED452C to my broader Puzzle Plan and EQUITAS work. My research interests are not separate from my teaching interests. In both places, I am asking how systems define competence, whose evidence gets recognized, and how design can reduce inequity before harm occurs. In education, that means designing classrooms where neurodivergent, culturally and linguistically diverse, disabled, chronically ill, multilingual, and multiply identified students are not treated as exceptions to be managed. In my larger AI and equity work, it means designing systems that do not mistake difference for deficiency. ED452C gave me a practical teaching language for the same ethical commitment that drives my research. Going forward, I want to continue reading about UDL, self-determination, inclusive assessment, transition planning, culturally sustaining pedagogy, and SEL as classroom-climate design. I also want to keep refining practical tools: agency audits, evidence-plan checklists, AI prompt routines, student voice prompts, and reflection structures that help students name how they learn. My next step is not simply to know more theory. It is to turn the theory into repeatable classroom habits.

Ultimately, what I gained from ED452C is a clearer professional stance. Inclusion is not proven by where a student sits. It is proven by what the environment allows the student to access, decide, express, revise, and become. If I carry that stance into my future teaching, then I will be better prepared to design classrooms where students are not merely present, but genuinely recognized as capable participants in their own learning.

\reviewsection{References}
\small
Rivers, S. E., Brackett, M. A., Reyes, M. R., Elbertson, N. A., \& Salovey, P. (2013). Improving the social and emotional climate of classrooms: A clustered randomized controlled trial testing the RULER approach. \textit{Prevention Science, 14}(1), 77--87. https://doi.org/10.1007/s11121-012-0305-2

Wehmeyer, M. L., \& Kurth, J. A. (2021). \textit{Inclusive education in a strengths-based era: Mapping the future of the field}. Norton.

Wehmeyer, M. L., \& Shogren, K. A. (2013). Self-determination and learners with autism spectrum disorder. Course reading for ED452C.

\end{document}
